\documentclass[11pt,a4paper]{article}

\usepackage[utf8]{inputenc}
\usepackage[T1]{fontenc}
\usepackage[margin=1in]{geometry}
\usepackage{amsmath}
\usepackage{booktabs}
\usepackage{hyperref}
\usepackage{tabularx}

\hypersetup{
  colorlinks=true,
  linkcolor=blue,
  citecolor=blue,
  urlcolor=blue
}

\title{\textbf{FileRerankingBench:\\A Small Benchmark for File Selection in Code-Editing Agents}}
\author{Ahmad Jiha \and Daniel Chen\\Anything, Inc.}
\date{}

\begin{document}
\maketitle

\begin{abstract}
Code-editing LLM agents need to select a small set of relevant project files
before editing. FileRerankingBench is a small public benchmark for that
file-selection step. The benchmark contains prompt-level relevance labels,
candidate file metadata, a project corpus, and aggregate evaluation numbers.
The public artifact is intentionally limited to benchmark data, aggregate
results, and this paper.
\end{abstract}

\section{Motivation}

Code-editing agents often operate on projects that are too large to place
entirely in model context. Before the model edits code, a system must decide
which files are likely to matter for the user's request. Poor file selection can
cause the agent to miss the relevant file, edit the wrong file, or waste context
on unrelated boilerplate.

FileRerankingBench isolates this pre-edit file-selection problem. Each benchmark
case provides a user prompt, a set of candidate files with timestamp metadata,
and the expected relevant file paths.

\section{Dataset}

The current public dataset contains one small hiking-app project and five query
fixtures. The cases cover:

\begin{itemize}
\item direct lexical matches between the prompt and target page,
\item stale or misleading recency metadata,
\item paraphrased references to UI text,
\item and a prompt whose requested page does not exist in the candidate set.
\end{itemize}

The fixture file is \texttt{data/file-reranking-evals.json}. The project corpus
is \texttt{data/projects/hiking-app-small.zip}. Candidate paths in the fixture
are app-root paths, for example
\texttt{/apps/web/pages/index/+Page.jsx}; the zip stores files below
\texttt{createxyz-project/}.

\section{Metrics}

The benchmark reports Recall@$k$ and Precision@$k$ for
$k \in \{1, 3, 5, 10\}$. Let $T_k$ be the top-$k$ returned files and let $R$ be
the set of expected relevant files.

\[
\mathrm{Recall}@k =
\begin{cases}
1 & \text{if } |R| = 0,\\
\frac{|T_k \cap R|}{|R|} & \text{otherwise.}
\end{cases}
\]

\[
\mathrm{Precision}@k = \frac{|T_k \cap R|}{|T_k|}
\]

If fewer than $k$ candidate files are available, $T_k$ contains only the
available candidates.

\section{Published Results}

Table~\ref{tab:results} reports aggregate reference numbers over the current
public fixture set. These numbers are intended to anchor the benchmark and make
future comparisons explicit without expanding the public artifact beyond data,
results, and this paper.

\begin{table}[h]
\centering
\small
\begin{tabularx}{\textwidth}{lrrrr}
\toprule
System & Recall@1 & Precision@1 & Recall@3 & Precision@3 \\
\midrule
Original candidate order & 0.6000 & 0.4000 & 1.0000 & 0.2667 \\
Recency baseline & 0.4000 & 0.2000 & 1.0000 & 0.2667 \\
Lexical baseline & 0.6000 & 0.4000 & 1.0000 & 0.2667 \\
Rank-all oracle upper bound & 1.0000 & 0.8000 & 1.0000 & 0.2667 \\
\bottomrule
\end{tabularx}
\caption{Macro-averaged results over five public benchmark cases. Recall@5,
Precision@5, Recall@10, and Precision@10 match the Recall@3 and Precision@3
columns in this corpus because each case has three candidate files.}
\label{tab:results}
\end{table}

\section{Release Contents}

The public repository at
\url{https://github.com/Create-Inc/file-reranking-bench} contains:

\begin{itemize}
\item the benchmark fixture JSON,
\item the project corpus zip,
\item aggregate result tables,
\item and this paper.
\end{itemize}

The release does not include implementation source code or runtime details.

\section{Limitations}

The current corpus is intentionally small: five queries over one project. It is
best understood as a regression benchmark and publication artifact, not a broad
statistical measurement of file-selection quality. Future versions should add
more projects, more user prompts, multi-file relevance labels, and richer
negative cases.

\section{Conclusion}

FileRerankingBench provides a compact public artifact for evaluating the
file-selection step in code-editing agents. By publishing the data and aggregate
results, the benchmark can support external comparison while keeping the release
focused on reproducible benchmark artifacts.

\end{document}
